% Generated draft: requires mathematical and source review.
\documentclass[11pt]{article}
\usepackage[margin=25mm]{geometry}
\usepackage{fontspec}
\setmainfont{Latin Modern Roman}
\usepackage{amsmath,amssymb,mathtools,unicode-math}
\setmathfont{Latin Modern Math}
\usepackage{enumitem,needspace,xurl,hyperref}
\setlength{\parindent}{0pt}
\setlength{\parskip}{5pt}
\setlength{\emergencystretch}{3em}
\newcommand{\N}{\mathbb{N}}
\newcommand{\Z}{\mathbb{Z}}
\newcommand{\Q}{\mathbb{Q}}
\newcommand{\R}{\mathbb{R}}
\newcommand{\C}{\mathbb{C}}
\newcommand{\F}{\mathbb{F}}
\newcommand{\A}{\mathbb{A}}
\newcommand{\PP}{\mathbb{P}}
\newcommand{\E}{\mathbb{E}}
\newcommand{\eps}{\varepsilon}
\newcommand{\dd}{\,\mathrm{d}}
\begin{document}
\section*{UnsolvedMath ID: 20000328}
\textbf{Source problem code:} catalogue entry
\subsection{Definitions and mathematical statement}
Viehweg asks whether the fundamental group of a moduli stack $\mathcal M_H$ of canonically polarized manifolds can be almost abelian, and more precisely conjectures a negative answer for a positive-dimensional nonsingular variety $U$ mapping generically finitely to that stack. A group is almost abelian if it contains an abelian subgroup of finite index. In the usual complex interpretation use topological fundamental groups, with the stack/orbifold group for $\mathcal M_H$. Interpret generically finite as finite generic fibers over the image, so a corresponding smooth projective family with ample fiberwise canonical class has maximal variation. This is different from requiring the coarse moduli map to be quasi-finite everywhere. The ground field, group convention, connectedness and exact image condition are not fixed in the original.

The original reports the conjecture known for $\mathcal M_g$ and under suitable local Torelli hypotheses. It discusses $U=\PP^n\setminus D$ for a normal-crossings divisor with fewer than $n$ components, attributing results to Viehweg--Zuo [VZ02], and cites [VZ05] as further motivation. Its repeated component-bound sentence uses $H$ where the boundary notation is unclear; it should not be interpreted as a new strict/non-strict range.

The attached unreviewed aid treats $n\ge2$ and a nonempty reduced simple-normal-crossings divisor $D=\sum_{i=1}^r D_i\subset\PP^n_\C$, with component degrees $d_i$. Attributing the group calculation to Libgober, it gives
\[
\pi_1(U)\simeq\Z^r/\langle(d_1,\ldots,d_r)\rangle
\simeq\Z^{r-1}\oplus\Z/\gcd(d_1,\ldots,d_r)\Z.
\]
Thus these complements have abelian fundamental group. If such $U$ supports a maximal-variation family of canonically polarized manifolds, the aid invokes Viehweg--Zuo and log-general-type results to require
\[
r\ge n,\qquad\sum_i d_i\ge n+2.
\]
For general-position hyperplanes, under the stronger everywhere quasi-finite coarse-moduli-map hypothesis, it reports Brody hyperbolicity and an entire-curve argument forcing $r\ge2n+1$. At the first allowed hyperplane number $r=2n+1$, its group formula yields $\pi_1(U)\simeq\Z^{2n}$. These are necessary conditions for a possible example, not a construction of a family or a counterexample to the conjecture. The all-variety question and the stack-group question remain separate from this boundary calculation.
\subsection{Short English statement}
Ask whether a positive-dimensional maximal-variation base for canonically polarized manifolds can have virtually abelian fundamental group. The aid gives necessary boundary conditions for projective-space complements and a stronger hyperplane bound under an everywhere quasi-finite moduli map.
\subsection{Sources}
\begin{itemize}
\item Viehweg, Problem 2.1 in AIM \emph{Compact moduli spaces and birational geometry}, Section 2, ``Moduli spaces of manifolds'': \url{https://aimath.org/WWN/birational/birational.pdf}. Canonical location: \texttt{aim-algebraic-geometry-notes.json}, notes[327]. Supplied verification resolves the split problem number; no PDF was consulted.
\item Original references: Viehweg and Zuo [VZ02] and [VZ05], without complete bibliographic entries. Original claims concerning $\mathcal M_g$, local Torelli and the boundary component restriction are retained as source reports.
\item The attached aid invokes Libgober, Viehweg--Zuo and log-general-type results without full citations, and proposes the combined necessary conditions with low novelty confidence. It is machine generated and not peer reviewed.
\item Local registry: \texttt{data/unsolvedmath/problems.json}, permanent catalogue entry, exact source code \texttt{catalogue entry}.
\end{itemize}
\end{document}
